% =====================================================================
%  Etnografi, Demografi, dan Profil Subjek Lapangan
%  Instrumen kerja lapangan, Koridor Jalan Jenderal Sudirman
%  Apple Developer Academy @ BINUS, Cohort 9 - Tim Riset Pijak
%
%  Target lapangan: Jumat, 19 Juni 2026.
%  Dokumen ini menyatukan tiga catatan kerja tim:
%    - Atribut Subjek (grid kuota/screening)
%    - Etnografi Subjek (profil persona, siapa ditemui & apa ditanya)
%    - Protokol Fieldcraft (penampilan, pendekatan, psikologi intercept)
%
%  Kompilasi di OVERLEAF (pdfLaTeX). Tanpa Biber: sumber dikutip inline.
%  babel option 'indonesian' (BUKAN 'bahasa').
% =====================================================================
\documentclass[11pt,a4paper]{report}

\usepackage[utf8]{inputenc}
\usepackage[T1]{fontenc}
\usepackage[indonesian]{babel}
\usepackage[margin=2.2cm]{geometry}   % beri ruang untuk tabel lebar
\usepackage{graphicx}
\usepackage{array}
\usepackage{tabularx}
\usepackage{longtable}
\usepackage{booktabs}
\usepackage{enumitem}
\usepackage{fancyhdr}
\usepackage[hidelinks]{hyperref}

\setlist{itemsep=2pt, topsep=3pt}
\setlength{\tabcolsep}{4pt}           % rapatkan kolom agar matriks muat
\renewcommand{\arraystretch}{1.1}

% --- Running headers/footers (sejajar dengan main_id.tex) ----------------
\newcommand{\shorttitle}{Profil Subjek Lapangan, Koridor Sudirman}
\newcommand{\teamname}{Tim Riset Pijak}
\pagestyle{fancy}
\fancyhf{}
\fancyhead[L]{\footnotesize \shorttitle}
\fancyhead[R]{\footnotesize \teamname}
\fancyfoot[C]{\footnotesize \thepage}
\renewcommand{\headrulewidth}{0.4pt}
\fancypagestyle{plain}{\fancyhf{}\fancyfoot[C]{\footnotesize \thepage}%
  \renewcommand{\headrulewidth}{0pt}}

% Kartu persona ringkas: label tebal + isi
\newcommand{\fld}[1]{\smallskip\noindent\textbf{#1.}\ }

\begin{document}
\pagenumbering{roman}

% --- Halaman judul -------------------------------------------------------
\begin{titlepage}
\thispagestyle{empty}
\centering
\vspace*{2cm}

{\Large\bfseries Etnografi, Demografi, dan Profil Subjek Lapangan\par}
\vspace{0.5em}
{\large Koridor Jalan Jenderal Sudirman dan Jalan-Jalan di Belakangnya\par}
\vspace{0.3em}
{\normalsize\itshape Siapa yang ditemui, apa yang ditanya, dan apa yang diamati\par}
\vspace{1.6cm}

{\large Urban Innovation Challenge (Challenge 4)\par}
{\normalsize Challenge Based Learning, Instrumen Kerja Lapangan\par}
\vspace{0.6em}
{\normalsize Target lapangan: Jumat, 19 Juni 2026\par}
\vspace{2cm}

{\large\bfseries \teamname\par}
\vspace{0.8em}
{\normalsize
Gading Aditya Perdana \\
Kenrich Heavenly Sandria \\
Naufal Ammar Rauf \\
Jatayu Muhammad Wicaksono \\
Sintani Gina Lestari \par}

\vfill
{\footnotesize Dokumen pendamping tinjauan pustaka dan metodologi (Bab II--III).
Pertanyaan ditulis netral dan dua arah; asumsi disimpan pada lembar terpisah,
tidak pernah di dalam pertanyaan.\par}
\end{titlepage}

\tableofcontents
\clearpage
\pagenumbering{arabic}

% =====================================================================
\chapter{Tujuan dan Prinsip Penyampelan}

\section{Tujuan Dokumen}
Dokumen ini memetakan ekosistem manusia di koridor Sudirman agar setiap anggota
tim tahu, untuk tiap persona: di mana dan kapan mereka benar-benar muncul, apakah
mereka di-\emph{wawancara} atau di-\emph{observasi}, kalimat pembuka yang dipakai
(dalam Bahasa Indonesia), serta sensitivitas yang harus dijaga. Dokumen menyatukan
tiga catatan kerja tim: grid atribut/kuota, profil persona, dan protokol \emph{fieldcraft}.

\section{Dua Fakta yang Mengatur Seluruh Penyampelan}
Kesenjangan literatur kami adalah soal \textbf{gradien} (trotoar lebar Sudirman vs
jalan satu blok di belakangnya), dan koridor berjalan menurut \textbf{jam}: siapa
yang hadir berubah total tiap waktu. Menyampel tempat yang salah pada jam yang salah
menghasilkan populasi yang tidak ada di waktu lain. Bab~2 menurunkan kedua fakta ini
menjadi peta ruang dan peta waktu.

\section{Dua Sumbu yang Menjadi Sel Kuota}
Seluruh atribut runtuh menjadi dua sumbu yang benar-benar memisahkan populasi kami,
dan keduanya tidak pernah disilangkan oleh literatur terkumpul:
\begin{enumerate}
  \item \textbf{Berjalan karena pilihan vs berjalan karena keharusan.} Pelari Car
    Free Day dan komuter hari kerja menapak aspal yang sama, tetapi subjek yang
    berlawanan. Tandai setiap responden.
  \item \textbf{Panggung depan vs panggung belakang.} Komuter sumbu utama vs warga
    Karet Tengsin adalah keseluruhan gradien kami. Kelas, asal, dan kualitas
    infrastruktur ikut bergerak sepanjang sumbu ini.
\end{enumerate}
Silangkan keduanya dengan \textbf{gender} (yang kami \emph{oversample}) dan
\textbf{jendela waktu}, dan terbentuk grid kuota yang dapat dipertanggungjawabkan,
mengisi tepat kesenjangan yang ditunjuk Bab~II.

% =====================================================================
\chapter{Lanskap Koridor: Gradien Ruang dan Jam}

\section{Gradien Ruang (Panggung Depan ke Panggung Belakang)}
Cerita yang harus diuji data: sumbu utama dirancang untuk pejalan kaki, panggung
belakang \emph{dihuni} oleh mereka, dan keduanya menyatu secara ekonomi. Kampung
bertahan dari menara; pekerja menara makan dan parkir di kampung. Berjalan menyusuri
\emph{seam} di antara keduanya berarti menyusuri pertanyaan penelitian kami.

\begin{longtable}{@{}p{2.8cm}p{4.2cm}p{3.2cm}p{4.3cm}@{}}
\toprule
\textbf{Zona} & \textbf{Apa itu} & \textbf{Siapa di sini} & \textbf{Mengapa penting} \\
\midrule
\endfirsthead
\toprule
\textbf{Zona} & \textbf{Apa itu} & \textbf{Siapa di sini} & \textbf{Mengapa penting} \\
\midrule
\endhead
\textbf{Sumbu utama} (trotoar Jl. Jend. Sudirman) &
Trotoar revitalisasi 8--12 m, \emph{bollard}, ``panggung depan'' Jakarta &
Pekerja kantoran transit, pelari, kota yang difoto &
Basis ``baik''. Terlihat \emph{walkable}; tanyakan apakah \emph{terasa} \emph{walkable}. \\
\addlinespace
\textbf{Simpul transit} (Dukuh Atas / Stasiun Sudirman / Jl.\ Blora / Jl.\ Kendal) &
Tempat KRL, MRT, LRT, TransJakarta, dan ojol bertabrakan &
Komuter \emph{last-mile}, pengemudi ojol, pedagang &
Friksi dan densitas tertinggi. Tumpukan ``menunggu ojol'' (Kompas, Mei 2025) terjadi
di sini, meluber ke Jl.\ Blora. \\
\addlinespace
\textbf{Seam} (Benhil / Karet, sekitar Halte Bendungan Hilir, JPO Benhil) &
Tempat sumbu mengkilap bertemu perdagangan lama &
PKL, warteg, ekosistem Pasar Benhil &
Sabuk PKL malam menempati 100--200 m trotoar di bawah lampu lapak, ``perpindahan''
yang diramalkan tinjauan pustaka. \\
\addlinespace
\textbf{Panggung belakang} (Karet Tengsin, gang di belakang Sudirman Park, tepi Kali Krukut) &
Kampung padat tepat satu tembok di belakang menara &
Warga kampung, juru parkir liar, penghuni kosan &
Gradien kami yang kurang terukur. ``Dua wajah dipisah satu tembok'': parkir motor
Rp10k/hari, warteg menyuapi pekerja kantor, tekanan gentrifikasi, tepi sungai rawan
banjir (Mojok, Feb 2026; Kumparan, Apr 2025). \\
\bottomrule
\end{longtable}

\section{Jam (Segmen yang Sama adalah Lima Tempat Berbeda)}

\begin{longtable}{@{}p{3.4cm}p{4.2cm}p{6.4cm}@{}}
\toprule
\textbf{Jendela} & \textbf{Populasi dominan} & \textbf{Perilaku khas untuk ditangkap} \\
\midrule
\endfirsthead
\toprule
\textbf{Jendela} & \textbf{Populasi dominan} & \textbf{Perilaku khas untuk ditangkap} \\
\midrule
\endhead
\textbf{Puncak pagi} 06.30--09.00 (puncak \textasciitilde07.30) &
Komuter kantoran turun KRL; cepat, ber-\emph{headphone}, terarah &
Arus KRL$\rightarrow$kantor dan KRL$\rightarrow$MRT; jembatan penghubung Dukuh Atas
``mengular''; rebutan jemput ojol di Jl.\ Blora \\
\addlinespace
\textbf{Tengah hari panas} 11.00--14.00 &
Pekerja makan siang, dagang warteg/PKL &
Mencari teduh, payung untuk panas, menyeberang ke sisi teduh, mundur ke
\emph{indoor}/\emph{skywalk} \\
\addlinespace
\textbf{Puncak sore} 15.00--18.30 &
Arus balik MRT$\rightarrow$KRL; PKL menggelar &
Lelah, melambat; sabuk PKL menyala dekat Benhil; Satpol PP hadir 17.00--21.00 \\
\addlinespace
\textbf{Malam} setelah 18.30 &
Lebih tipis, lebih maskulin, pelanggan PKL, warga &
Celah pencahayaan, divergensi rasa aman, perubahan rute perempuan \\
\addlinespace
\textbf{Akhir pekan / Car Free Day} (Min 06.00--10.00/11.00) &
Pelari, pesepeda, keluarga, komunitas, pasar Minggu, 30--90 ribu orang &
Koridor yang \emph{berbeda total}: rekreatif, sosial, padat. Baik untuk volume dan
\emph{intercept} komunitas, tidak mewakili realitas hari kerja \\
\bottomrule
\end{longtable}

\noindent\textbf{Peringatan tetap.} Hari WFH dan hari demonstrasi mengubah segalanya.
Detik (Sep 2025) mendokumentasikan Stasiun Sudirman nyaris kosong pada pagi berimbauan
WFH. Catat hari, cuaca, dan tiap sebutan ``kantor lagi WFH?'' pada setiap lembar
\emph{tally}, agar hitungan tetap dapat dibandingkan.

% =====================================================================
\chapter{Atribut Subjek dan Sel Kuota}

\section{Matriks Atribut Utama}
Tag berikut dipasang pada setiap responden agar sampel \emph{nonprobability} tetap
seimbang lintas persona, kelas, asal, moda, dan waktu.

\begin{longtable}{@{}p{2.6cm}p{1.7cm}p{1.6cm}p{2.4cm}p{2.3cm}p{2.1cm}@{}}
\toprule
\textbf{Persona} & \textbf{Usia} & \textbf{Gender} & \textbf{Kelas/proksi} & \textbf{Asal} & \textbf{Alasan jalan} \\
\midrule
\endfirsthead
\toprule
\textbf{Persona} & \textbf{Usia} & \textbf{Gender} & \textbf{Kelas/proksi} & \textbf{Asal} & \textbf{Alasan jalan} \\
\midrule
\endhead
Pekerja kantoran / komuter & 23--45 & Campuran & Kerah putih; menengah--atas & Komuter Bodetabek & Keharusan (\emph{last-mile}) \\
\addlinespace
Pejalan kaki perempuan \emph{(oversample)} & 20--45 & Perempuan & Kerah putih campuran & Komuter + lokal & Keharusan; rute dibentuk rasa aman \\
\addlinespace
Pengguna \emph{last-mile} transit & 22--45 & Campuran & Kerah putih menengah & Komuter Bodetabek & Keharusan \\
\addlinespace
Pengemudi ojol & 20--50 & Mayoritas pria & Sektor informal; bawah--menengah & Jabodetabek & Mereka \emph{adalah} \emph{last-mile} \\
\addlinespace
PKL (pedagang) & 30--60 & Campuran & Sektor informal; bawah & Sering eks Pasar Benhil, lokal & Bekerja di \emph{seam} \\
\addlinespace
Warga kampung (Karet Tengsin) & Semua & Campuran & Bawah; penyewa kos/kontrak & Lokal Jakarta, menetap & Keharusan; infrastruktur terburuk \\
\addlinespace
Mobilitas khusus (disabilitas, lansia, kereta bayi) & Semua & Semua & Campuran & Campuran & Keharusan; uji \emph{walkability} paling jujur \\
\addlinespace
Pelajar / rekreasi CFD & 17--30 (hari kerja); semua (CFD) & Campuran & Campuran & Campuran & \textbf{Pilihan} (rekreasi) \\
\addlinespace
Satpam / porter / staf pendukung & 25--55 & Mayoritas pria & Jasa; bawah--menengah & Mayoritas lokal & Diam; saksi koridor \\
\bottomrule
\end{longtable}

\section{Kehadiran menurut Waktu dan Zona}
\emph{$\bullet$ = hadir, $\bullet\bullet\bullet$ = dominan, $\circ$ = tipis/absen.}

\begin{longtable}{@{}p{3.2cm}ccccc p{2.6cm}@{}}
\toprule
\textbf{Persona} & \textbf{Pagi} & \textbf{Siang} & \textbf{Sore} & \textbf{Malam} & \textbf{CFD} & \textbf{Zona utama} \\
\midrule
\endhead
Pekerja kantoran & $\bullet\bullet\bullet$ & $\bullet\bullet$ & $\bullet\bullet\bullet$ & $\circ$ & $\circ$ & Sumbu + simpul \\
Pejalan kaki perempuan & $\bullet\bullet\bullet$ & $\bullet\bullet$ & $\bullet\bullet\bullet$ & $\bullet$ \emph{(kunci)} & $\bullet$ & Sumbu + jalan belakang \\
\emph{Last-mile} transit & $\bullet\bullet\bullet$ & $\bullet$ & $\bullet\bullet\bullet$ & $\circ$ & $\circ$ & Dukuh Atas / Jl.\ Blora \\
Ojol & $\bullet\bullet\bullet$ & $\bullet\bullet$ & $\bullet\bullet\bullet$ & $\bullet\bullet$ & $\bullet$ & Mulut stasiun, Jl.\ Blora \\
PKL & $\circ$ & $\bullet$ & $\bullet\bullet\bullet$ & $\bullet\bullet\bullet$ & $\bullet\bullet$ & \emph{Seam} Benhil / Karet \\
Warga kampung & $\bullet$ & $\bullet\bullet$ & $\bullet\bullet$ & $\bullet\bullet\bullet$ & $\bullet$ & Gang Karet Tengsin \\
Mobilitas khusus & $\circ$ & $\circ$ & $\circ$ & $\circ$ & $\bullet$ & Penyeberangan, \emph{ramp} \\
Pelajar / CFD & $\circ$ & $\bullet$ & $\bullet$ & $\bullet$ & $\bullet\bullet\bullet$ & Dukuh Atas, FX \\
Satpam / pendukung & $\bullet\bullet$ & $\bullet\bullet$ & $\bullet\bullet$ & $\bullet\bullet$ & $\bullet$ & Pintu gedung/stasiun \\
\bottomrule
\end{longtable}

\section{Penanda Lapangan (Atribut Teramati, Kenali Tanpa Bertanya)}
Atribut yang dapat dilihat langsung menentukan apakah subjek menjadi \emph{tally}
atau wawancara.
\begin{itemize}
  \item \textbf{Pekerja kantoran:} busana kantor (\emph{lanyard}/ID, blazer, \emph{business-casual}), ransel/tote, sibuk menavigasi ponsel, langkah cepat terarah, \emph{headphone}, JakLingko/uang elektronik di tangan.
  \item \textbf{Pejalan kaki perempuan:} tambahan, menyesuaikan rute memutari halangan, postur kunci/ponsel di tangan, perubahan langkah dan rute saat senja pada segmen yang sama.
  \item \textbf{Komuter terjebak:} berdiri di tepi trotoar, berulang mengecek aplikasi \emph{ride-hail}, berkumpul dekat mulut stasiun.
  \item \textbf{Ojol:} jaket hijau (Gojek)/hijau-putih (Grab)/lainnya, helm, parkir berjajar, mengamati \emph{ping} aplikasi.
  \item \textbf{PKL:} gerobak/lapak tetap, lampu terang setelah gelap, makanan/minuman/rokok, jejak 100--200 m di trotoar Benhil.
  \item \textbf{Warga kampung:} keluar-masuk gang, penjaga warteg, juru parkir menarik Rp10k, busana rumah santai.
  \item \textbf{Mobilitas khusus:} kursi roda, tongkat, \emph{walker}, kereta bayi; tampak menegosiasi \emph{ramp}, ubin pemandu, tinggi \emph{curb}.
  \item \textbf{Pejalan CFD:} pakaian olahraga, perlengkapan lari/sepeda, jersey komunitas, keluarga, hanya Minggu pagi.
  \item \textbf{Satpam / pendukung:} seragam, pos tetap, sedang bertugas.
\end{itemize}

\section{Atribut Laten (Hanya Diketahui dengan Bertanya, Muatan \emph{Intercept})}
\begin{longtable}{@{}p{4.2cm}p{4.0cm}p{5.0cm}@{}}
\toprule
\textbf{Atribut} & \textbf{Paling penting bagi} & \textbf{Cara mengetahui} \\
\midrule
\endhead
Asal sebenarnya \& jalur \emph{last-mile} & Pekerja kantoran, pengguna transit & ``Tadi jalan dari mana ke mana?'' \\
Jarak toleransi berjalan & Semua; perempuan (lebih pendek) & ``Tadinya mau jalan sampai mana sebelum nyerah / naik ojol?'' \\
Logika memilih \& menghindari rute & Perempuan (rasa aman), semua & ``Ada bagian rute yang dipilih atau dihindari? Kenapa?'' \\
Perubahan rute siang vs malam & Perempuan & Segmen sama, dua jendela \\
Perilaku navigasi & Pekerja kantoran & ``Pakai aplikasi atau udah hafal?'' \\
Mengatasi panas / hujan & Semua & Amati payung/perpindahan teduh, lalu tanya \\
Simbiosis ekonomi & Kampung, PKL & ``Pelanggan kebanyakan orang gedung atau warga sini?'' \\
Pengalaman penertiban & PKL & Halus, ``biasanya gelar jam berapa sampai jam berapa?'' \\
Tekanan lahan / penggusuran & Kampung & Hanya jika \emph{mereka} yang mengangkat \\
\bottomrule
\end{longtable}

% =====================================================================
\chapter{Profil Subjek per Persona}

Untuk tiap persona: profil, lokasi/waktu, metode (\emph{wawancara} = \emph{intercept};
\emph{observasi} = hitung/audit tanpa kontak), pertanyaan, dan sensitivitas.
Pertanyaan ditulis netral dan dua arah, tanpa asumsi.

\section{Pekerja Kantoran / Komuter \textnormal{(sampel tulang punggung)}}
\fld{Profil} Populasi penentu. Sudirman adalah submarket perkantoran terbesar Jakarta
(sekitar 2,9 juta m\textsuperscript{2}). Sebagian besar bukan warga lokal: naik KRL dari
Bogor, Bekasi, Depok, Tangerang, transit di Manggarai, turun di Stasiun Sudirman, lalu
berjalan atau melompat MRT/TransJakarta untuk ruas terakhir. Berjalan di sini adalah
penghubung \emph{last-mile}, bukan moda pilihan. Cenderung 23--45 tahun, berbusana
kantor, terburu waktu, menavigasi ponsel.

\fld{Lokasi \& waktu} Pintu keluar Stasiun Sudirman, jembatan penghubung Dukuh Atas, dan
trotoar utama, 06.30--09.00 dan 16.30--18.30. Mereka tidak berhenti lama; tangkap pada
jeda alami (menunggu ojol, di penyeberangan, di gerobak kopi, di peron).

\fld{Metode} Keduanya. Observasi arus pagi untuk \emph{tally} (trotoar vs badan jalan,
arah, gender, langkah). \emph{Intercept} mereka yang \emph{sudah menunggu}, itulah satu-
satunya jendela wawancara yang realistis.

\fld{Pertanyaan ($\leq$5 menit, setelah persetujuan lisan)}
\begin{itemize}
  \item ``Tadi Bapak/Ibu jalan kakinya dari mana ke mana?''
  \item ``Bagian mana dari rute tadi yang paling enak dilewati, dan bagian mana yang bikin males?''
  \item ``Kalau lagi buru-buru, jalurnya berubah nggak? Berubahnya ke mana?''
  \item ``Tadi pakai aplikasi buat nentuin jalan, atau udah hafal? Kalau pakai, yang mana?''
  \item ``Pernah nggak milih naik ojol buat jarak yang sebenernya deket? Kira-kira kenapa?''
\end{itemize}

\fld{Sensitivitas} Mereka terburu dan bisa mengira tim sebagai sales/MLM. Buka dengan
``kami mahasiswa lagi riset, boleh minta waktu sebentar?'' dan biarkan mereka tetap
berjalan bila perlu.

\section{Pejalan Kaki Perempuan \textnormal{(oversample; sumbu rasa aman)}}
\fld{Profil} Metodologi sudah menandai kelompok ini untuk \emph{oversample}, dan data
mendukung: pada survei KRPA empat dari lima perempuan melaporkan pelecehan di ruang
publik, sekitar 46,8\% di transportasi umum; jalanan, permukiman, dan transit adalah
tiga lokasi teratas (KRPA 2022). Norma berjalan Jakarta menunjukkan jarak terima
perempuan menuju stasiun lebih pendek (sekitar 593 m vs 629 m laki-laki); rasa aman dan
kenyamanan secara harfiah mengecilkan peta mereka (Afkara \& Kusuma 2020).

\fld{Lokasi \& waktu} Di mana pun kohort kantoran berada, ditambah satu \emph{pass}
\textbf{malam} dan perbandingan \textbf{siang vs malam} pada segmen yang \emph{sama}.
Kontras antara sumbu utama yang terang dan jalan belakang yang gelap adalah temuannya.

\fld{Metode} Keduanya, dengan wawancara \textbf{dipimpin Sintani}, berpasangan, terutama
pada malam dan topik sensitif.

\fld{Pertanyaan (netral, jangan menggiring, jangan sebut ``pelecehan'' lebih dulu)}
\begin{itemize}
  \item ``Kalau jalan sendiri di sini, ada nggak bagian rute yang Mbak/Ibu pilih atau justru hindari? Kenapa?''
  \item ``Beda nggak rasanya jalan di sini siang sama malam?''
  \item ``Hal apa yang bikin satu jalur kerasa lebih nyaman buat dilewatin?''
  \item \emph{(hanya jika mereka mengangkat)} ``Tadi sempat nyebut soal merasa kurang nyaman, boleh cerita lebih lanjut kalau berkenan?''
\end{itemize}

\fld{Sensitivitas} Kelompok paling hati-hati. Tanpa asumsi, tanpa menggali trauma, jalan
keluar mudah, perekaman hanya dengan persetujuan terpisah, wajah dikaburkan. Bila
responden tampak terganggu, hentikan dan ucapkan terima kasih.

\section{Pengguna \emph{Last-Mile} Transit \textnormal{(beririsan dengan kantoran, memuat momen ``terjebak'')}}
\fld{Profil} Kerumunan transit di simpul Dukuh Atas. Titik nyeri tanda tangan
terdokumentasi: penumpang menumpuk di tepi trotoar menunggu ojol yang tak kunjung datang,
sementara pengemudi berjaket hijau-kuning menjajari Jl.\ Blora (Kompas, Mei 2025).

\fld{Lokasi \& waktu} Pintu utara Stasiun Sudirman, Jl.\ Blora, mulut MRT Dukuh Atas,
terowongan penghubung. Puncak pagi dan sore.

\fld{Metode} Observasi tumpukan dan perilaku antre untuk \emph{tally}; \emph{intercept}
orang yang \emph{berdiri diam menunggu} (audiens paling tertahan).

\fld{Pertanyaan}
\begin{itemize}
  \item ``Lagi nunggu ojol, atau mau lanjut jalan/naik apa?''
  \item ``Dari turun kereta sampai ke kantor, biasanya paling ribet di bagian mana?''
  \item ``Kalau ojolnya susah didapet kayak gini, biasanya Bapak/Ibu ngapain?''
\end{itemize}

\fld{Sensitivitas} Rendah. Orang yang menunggu sering bersedia bicara untuk mengisi waktu.

\section{Ojol \textnormal{(aktor informal; infrastruktur dalam wujud manusia)}}
\fld{Profil} Mereka \emph{adalah} sistem \emph{last-mile} yang gagal disediakan trotoar.
Berkumpul di titik jemput yang dapat diprediksi (Jl.\ Blora, mulut stasiun), paling
paham mikro-geografi, dan punya bacaan struktural soal di mana pejalan kaki menyerah.
Catatan: ojol punya budaya buruh/aksi yang aktif di Jakarta, bersikap netral, jangan
ekstraktif.

\fld{Lokasi \& waktu} Klaster jemput dekat Stasiun Sudirman dan Dukuh Atas, sepanjang
hari, memuncak seiring pasang-surut komuter.

\fld{Metode} Keduanya. Mereka \emph{key informant} unggul soal perilaku pejalan kaki
(``orang nyerah jalan di titik mana'').

\fld{Pertanyaan}
\begin{itemize}
  \item ``Penumpang paling banyak nge-pick up di titik mana sini, Bang?''
  \item ``Banyak nggak yang sebenernya deket tapi tetep milih ojol? Menurut Abang kenapa?''
  \item ``Jam berapa paling rame orang nungguin ojol di sini?''
\end{itemize}

\fld{Sensitivitas} Jangan ganggu saat sedang menjemput. Dekati saat senggang. Hindari
kesan penegakan platform/polisi.

\section{PKL \textnormal{(aktor informal; cerita perpindahan)}}
\fld{Profil} Kelompok paling sarat etnografi. Sumbu utama Sudirman resmi dinyatakan
terlarang bagi pedagang pada 2018, dengan relokasi ``ke jalan-jalan di belakang'',
tepat gradien yang kami kaji. Banyak yang berakar pada pembersihan Pasar Benhil;
status mereka resmi ``liar'', dan Satpol PP berpatroli di ruas Benhil sekitar
17.00--21.00 (Tempo/Media Indonesia/Detik 2018, berlanjut). Lapak malam mereka, terang,
menempati 100--200 m trotoar dekat JPO Benhil. Mereka sekaligus ``halangan'' dalam
audit sekaligus sumber pencahayaan, mata yang mengawasi jalan, dan kehidupan. Pegang
kedua kebenaran itu.

\fld{Lokasi \& waktu} \emph{Seam} Benhil/Karet, sore menjelang malam. Waktu menggelar
(sekitar 16.00--18.00) paling baik untuk bicara, mereka belum diserbu pembeli.

\fld{Metode} Keduanya. Observasi jejak, pencahayaan, dan pengalihan pejalan kaki untuk
audit. Wawancara halus sebagai \emph{key informant} soal ritme koridor dan penertiban.

\fld{Pertanyaan (hormati kerentanan, jangan interogasi legalitas)}
\begin{itemize}
  \item ``Biasanya mulai gelar dagangan jam berapa, sampai jam berapa, Pak/Bu?''
  \item ``Rame-nya pembeli pas jam berapa? Orang kantoran atau warga sini?''
  \item ``Menurut Bapak/Ibu, orang yang lewat sini lebih enak jalannya pas ada lapak atau pas kosong?''
\end{itemize}

\fld{Sensitivitas} Tinggi. Nafkah mereka rapuh secara hukum. Jangan pernah mengesankan
tim melapor; jangan memotret dengan cara yang dapat mengenali mereka ke penertiban.
Bingkai sebagai memahami jalan, bukan mengaudit kepatuhan.

\section{Warga Jalan Belakang (Karet Tengsin) \textnormal{(gradien kurang terukur, dalam wujud)}}
\fld{Profil} Orang yang tak pernah diukur literatur. Di gang belakang Sudirman Park,
warga menjalankan warteg yang menyuapi pekerja kantor, menyewakan kos, dan mengelola
parkir liar sekitar Rp10k/hari (vs basement menara yang mahal), simbiosis di mana
``menara adalah ekonomi panggung belakang'' (Mojok Feb 2026; Kumparan Apr 2025).
Permukiman menyentuh Kali Krukut, rawan banjir, menghadapi tekanan gentrifikasi dan
pembelian lahan. Mereka menapaki infrastruktur terburuk koridor setiap hari tanpa
pilihan moda.

\fld{Lokasi \& waktu} Gang Karet Tengsin, titik warteg dan parkir, menjelang siang dan
tengah hari (dagang warung), malam (di rumah). Masuk lewat warung: duduk, pesan, bicara.

\fld{Metode} Terutama \textbf{wawancara/nongkrong} (warung sebagai \emph{fieldsite}).
Audit kondisi trotoar jalan mereka sebagai paruh ``belakang'' tiap segmen berpasangan.

\fld{Pertanyaan}
\begin{itemize}
  \item ``Bapak/Ibu tiap hari jalan kakinya ke arah mana? Lewat gang yang mana?''
  \item ``Pelanggan warung kebanyakan orang gedung atau warga sini?''
  \item ``Kalau hujan, jalanan di sini gimana? Pernah kebanjiran?''
  \item ``Beda banget ya jalan di depan (Sudirman) sama di sini, menurut Bapak/Ibu kenapa bisa begitu?''
\end{itemize}

\fld{Sensitivitas} Sedang. Lahan/penggusuran adalah topik bermuatan; jangan membukanya
lebih dulu; bila diangkat, dengarkan, jangan menasihati. Tim adalah tamu di kampung,
bukan auditor.

\section{Kelompok Mobilitas Khusus \textnormal{(uji \emph{walkability} paling jujur)}}
\fld{Profil} Catatan tempel ``Disability Friendly?'' di papan asumsi menunjuk ke sini.
Ubin pemandu ada di sumbu revitalisasi tetapi rutin terhalang (motor, PKL, mobil
parkir). Kelompok ini menyingkap apakah infrastruktur bekerja bagi siapa pun yang bukan
orang dewasa tanpa hambatan.

\fld{Lokasi \& waktu} Jarang dan tak terduga di jalan; mungkin perlu rekrut lewat titik
aksesibilitas transit, atau amati pengguna kursi roda/tongkat/kereta bayi di
penyeberangan dan \emph{ramp}. Jangan paksakan kuota yang menghasilkan nol responden
nyata; catat jujur bila tidak terjangkau.

\fld{Metode} Observasi kelaikan \emph{ramp}, kontinuitas jalur pemandu, waktu
penyeberangan. Wawancara hanya dengan kehati-hatian ekstra dan manfaat jelas bagi mereka.

\fld{Pertanyaan}
\begin{itemize}
  \item ``Bagian mana dari jalur ini yang paling susah dilewati buat Bapak/Ibu?''
  \item ``Kalau ada halangan di trotoar, biasanya Bapak/Ibu ngakalin gimana?''
\end{itemize}

\fld{Sensitivitas} Jangan membantu tanpa bertanya; jangan memperlakukan sebagai inspirasi
atau tontonan. Penilaian mereka adalah keahlian, bukan kisah sedih.

\section{Pelajar/Mahasiswa \& Pejalan Rekreasi CFD \textnormal{(volume dan kontras ``jalan untuk senang'')}}
\fld{Profil} Hari kerja: relatif sedikit pelajar di sumbu kantor ini. Akhir pekan:
koridor berubah, Car Free Day (Min 06.00--10.00/11.00) menarik 30--90 ribu pelari,
pesepeda, keluarga, dan komunitas, dengan pasar Minggu dari FX hingga Benhil, Dukuh Atas,
sampai Sarinah. Inilah berjalan sebagai rekreasi, kebalikan dari kerja \emph{last-mile}
hari kerja, kasus kontras yang berguna, bukan pengganti data hari kerja.

\fld{Lokasi \& waktu} CFD Minggu pagi, simpul Dukuh Atas dan FX (titik kumpul komunitas
paling ramai).

\fld{Metode} Keduanya, bagus untuk \emph{intercept} bervolume tinggi dan untuk menyampel
makna ``nyaman'' yang \emph{diungkapkan} saat berjalan adalah pilihan, bukan keharusan.

\fld{Pertanyaan}
\begin{itemize}
  \item ``Sering ke sini pas weekday juga, atau cuma pas CFD?''
  \item ``Apa yang bikin jalan kaki di sini pas CFD kerasa beda?''
  \item ``Kalau hari biasa, masih betah jalan kaki di sini nggak? Kenapa?''
\end{itemize}

\fld{Sensitivitas} Rendah. Suasana santai. Ingat, kesan CFD \emph{bukan} realitas hari
kerja, beri label jelas pada data.

\section{Satpam, Porter, Petugas Kebersihan, Staf Pendukung \textnormal{(saksi yang diam)}}
\fld{Profil} Sering terabaikan, tetapi berdiri di satu titik sepanjang hari dan menonton
seluruh siklus koridor. Keamanan gedung, porter stasiun, penyapu. \emph{Key informant}
luar biasa soal pola (``jam segini selalu rame'', ``yang nyebrang sembarangan biasanya
di sini'').

\fld{Lokasi \& waktu} Pintu gedung, gerbang stasiun, sepanjang hari. Senggang pertengahan
pagi/siang.

\fld{Metode} Wawancara sebagai \emph{key informant}.

\fld{Pertanyaan}
\begin{itemize}
  \item ``Bapak tiap hari di sini, jam berapa paling rame orang jalan kaki?''
  \item ``Ada nggak titik yang orang sering nyebrang sembarangan atau jalan di badan jalan? Di mana?''
  \item ``Kalau hujan atau panas banget, orang-orang biasanya ngapain?''
\end{itemize}

\fld{Sensitivitas} Rendah, tetapi mereka mungkin sedang bertugas, jaga singkat, jangan
membuat mereka bermasalah dengan atasan.

\section{Matriks Utama (Siapa $\times$ Kapan $\times$ Di Mana $\times$ Metode)}
\begin{longtable}{@{}p{3.0cm}p{3.2cm}p{3.4cm}p{2.6cm}p{1.8cm}@{}}
\toprule
\textbf{Persona} & \textbf{Jendela terbaik} & \textbf{Simpul terbaik} & \textbf{Metode} & \textbf{Pemimpin} \\
\midrule
\endhead
Pekerja kantoran & 06.30--09.00; 16.30--18.30 & Stasiun Sudirman, jembatan, sumbu utama & Observasi + \emph{intercept} & Siapa pun \\
Pejalan kaki perempuan & Semua + \textbf{malam} + pasangan siang/malam & Sumbu + jalan belakang & \emph{Intercept} & \textbf{Sintani} (malam/sensitif) \\
\emph{Last-mile} transit & Puncak & Jl.\ Blora, mulut Dukuh Atas & Observasi + \emph{intercept} & Siapa pun \\
Ojol & Sepanjang hari, puncak & Jl.\ Blora, mulut stasiun & \emph{Key informant} & Siapa pun (senggang) \\
PKL & 16.00--21.00 & \emph{Seam} Benhil/Karet & Observasi + wawancara halus & Siapa pun (berpasangan) \\
Warga kampung & Siang \& malam & Gang Karet Tengsin, warteg & Nongkrong/wawancara & Berpasangan \\
Mobilitas khusus & Oportunistik & Penyeberangan, \emph{ramp}, transit & Observasi + wawancara hati-hati & Siapa pun \\
Pelajar / CFD & \textbf{Minggu 06.00--10.00} & Dukuh Atas, FX & Observasi + \emph{intercept} & Siapa pun \\
Satpam / staf & Pertengahan pagi/siang & Pintu gedung \& stasiun & \emph{Key informant} & Siapa pun \\
\bottomrule
\end{longtable}

% =====================================================================
\chapter{Observasi Tanpa Wawancara (Protokol Senyap)}
Sebagian data terbaik tidak butuh persetujuan dan tidak butuh percakapan. Saat satu
anggota mewawancara, anggota lain menjalankan \textbf{\emph{tally} 5 menit} dan
mengamati:
\begin{itemize}
  \item \textbf{Trotoar vs badan jalan:} berapa yang berjalan di badan jalan padahal
    trotoar tersedia, dan \emph{di mana persisnya} mereka meninggalkannya (halangan atau
    titik sempit yang mendorong mereka turun).
  \item \textbf{Perilaku menyeberang:} titik menyeberang sembarangan, jarak antar
    penyeberangan aman (tolok ukur 100--200 m), kepatuhan sinyal; relevan terhadap datum
    nasional bahwa mayoritas korban pejalan kaki terkait menyeberang di tengah ruas.
  \item \textbf{Serah-terima ojol:} di mana berjalan berakhir dan motor dimulai, dan
    sependek apa trip itu sebenarnya.
  \item \textbf{Mencari teduh:} perpindahan sisi tengah hari, penggunaan payung untuk
    panas, mundur ke \emph{skywalk}/\emph{indoor}, proksi teramati untuk ketidaknyamanan
    termal.
  \item \textbf{Inventaris halangan:} motor parkir di ubin pemandu, jejak PKL, \emph{ramp}
    rusak/terhalang, genangan air.
  \item \textbf{Spasi dan waktu bergender:} apakah jumlah perempuan turun setelah gelap
    pada satu segmen? Pada segmen mana, dan setajam apa?
  \item \textbf{Foto kontras depan/belakang:} potret jam yang sama, aktivitas yang sama,
    di sumbu dan satu blok di belakangnya. Citra berpasangan itulah temuan gradien.
\end{itemize}

% =====================================================================
\chapter{Strategi Pendekatan dan Wawancara}
\emph{Bingkai etis lebih dulu: teknik berikut adalah cara membangun \emph{rapport} dan
mengingat akurat, bukan tuas untuk mengorek pengakuan dari orang rentan. Dengan PKL dan
warga kampung khususnya, \emph{rapport} tidak boleh berubah jadi ekstraksi. Persetujuan
berlangsung terus; berhenti begitu mereka mau berhenti.}

\section{Sebelum Bicara: Pendekatan}
\begin{itemize}
  \item \textbf{\emph{Intercept} pada jeda alami, bukan saat melangkah.} Orang yang
    menunggu ojol, di penyeberangan, atau di warung punya perhatian luang; orang berjalan
    cepat tidak. Tumpukan di Stasiun Sudirman/Jl.\ Blora adalah audiens menganggur yang
    tertahan, titik hasil tertinggi.
  \item \textbf{Proksemik dan sudut tubuh.} Dekati dari depan agak menyerong, jangan dari
    belakang, jangan menghalangi jalan. Jaga sekitar 1 m, tangan terlihat, postur terbuka,
    senyum tulus kecil. Untuk responden perempuan ini penting: jangan mereproduksi dinamika
    menyudutkan yang membuat jalanan terasa tidak aman.
\end{itemize}

\section{Pembukaan dan Persetujuan (\emph{Foot-in-the-Door})}
\begin{itemize}
  \item \textbf{Mulai dengan permintaan kecil} (Cialdini): ``Permisi, boleh minta waktu
    2 menit?'' Satu ``ya'' kecil menyiapkan ``ya'' yang lebih besar.
  \item \textbf{Transparansi sebagai pemercepat kepercayaan.} Sebut dalam satu tarikan
    napas: siapa (mahasiswa, riset kampus Apple Developer Academy), apa (soal pengalaman
    jalan kaki di sini), berapa lama (sebentar), bahwa anonim, sukarela, dan bisa berhenti
    kapan saja.
  \item \textbf{Biarkan \emph{clipboard} menjelaskan tim.} Prop melegitimasi gangguan
    sebelum kalimat selesai.
\end{itemize}

\section{\emph{Rapport}: Mesin yang Sesungguhnya}
\begin{itemize}
  \item \textbf{Bangun tiga komponen} (Tickle-Degnen \& Rosenthal): perhatian (fokus
    penuh), positivitas (kehangatan), koordinasi (selaraskan ritme mereka).
  \item \textbf{Cocokkan register; alih kode.} Bahasa halus formal (Bapak/Ibu) dengan
    lansia, pekerja kantoran, pejabat; santai dengan ojol, PKL muda, pelajar. Terlalu
    formal terasa interogasi; terlalu santai dengan lansia terasa tidak hormat.
  \item \textbf{Pemanasan dengan obrolan tulus.} ``ramai ya pagi ini'', komentar soal
    panas, antrean, dagangan. Dengan pedagang, \textbf{beli dulu}: resiprositas plus
    hormat pada nafkah mereka.
  \item \textbf{\emph{Minimal encouragers} dan OARS} (Miller \& Rollnick): angguk, ``iya
    terus?'', ``menarik'', refleksi singkat. Bicara lebih sedikit dari mereka.
\end{itemize}

\section{Desain Pertanyaan dan \emph{Recall}}
\begin{itemize}
  \item \textbf{Corong (\emph{funnel}).} Buka luas, persempit kemudian: mulai ``ceritain
    dong tadi jalannya'', akhiri dengan hal spesifik.
  \item \textbf{Pertanyaan \emph{grand-tour}} (Spradley): minta mereka memandu satu
    kejadian nyata terkini, ``coba ceritain tadi pagi dari turun kereta sampai ke kantor''.
    Ini menghasilkan perilaku, bukan opini, tepat data ``dinyatakan'' kami.
  \item \textbf{Buang asumsi.} ``Kenapa jalan di badan jalan?'' sudah mengandaikan;
    sebagai ganti ``tadi lewat mana?'' dan biarkan alasannya muncul. Hindari pertanyaan
    menggiring, ganda, atau berjargon (``\emph{walkability}''), pakai kata mereka.
  \item \textbf{Reinstatemen konteks} (Fisher \& Geiselman): ``coba inget-inget pas keluar
    stasiun tadi, rame nggak, panas nggak?'' Bangun ulang adegan untuk akurasi ingatan.
    Tanyakan trip nyata \emph{terakhir}, bukan rata-rata.
\end{itemize}

\section{Topik Sensitif: Menurunkan Biaya Sosial Kejujuran}
\begin{itemize}
  \item \textbf{Bingkai orang ketiga/proyektif} (Tourangeau \& Yan, memangkas bias
    \emph{social-desirability}): ``ada yang bilang sini kurang nyaman buat perempuan kalau
    malam, menurut Mbak gimana?''
  \item \textbf{Normalisasi dulu:} ``pengalaman orang beda-beda soal ini.''
  \item \textbf{Rasa aman perempuan:} jangan sebut pelecehan lebih dulu; biarkan mereka
    memimpin. Bila diungkap, refleksikan dengan lembut dan jangan menggali trauma atau
    memperbesar distres. Tawarkan jalan keluar mudah dan ucapkan terima kasih.
  \item \textbf{Legalitas PKL/penggusuran kampung:} jangan pernah membingkai sebagai audit
    atau penegakan. Bila mereka mengangkat penggusuran, dengarkan, jangan menasihati, dan
    lindungi mereka, tanpa foto yang mengenali.
  \item \textbf{Lawan bias \emph{acquiescence}:} orang setuju demi sopan. Hindari giringan
    ya/tidak; tawarkan opsi seimbang, ``lebih enak yang mana, A, B, atau hal lain?''
\end{itemize}

\section{Menutup dan Koreografi Tim}
\begin{itemize}
  \item \textbf{Rangkum balik} (OARS): ``jadi intinya tadi, bener gitu, Pak?'', cek akurasi
    yang juga membuat mereka merasa didengar.
  \item \textbf{Buka lantai:} ``ada yang belum saya tanyain tapi penting menurut
    Bapak/Ibu?'', kutipan terbaik sering muncul di sini.
  \item \textbf{Berterima kasih, jangan janji berlebih.} Jangan mengesankan tim akan
    memperbaiki jalan mereka.
  \item \textbf{Berpasangan:} satu \emph{interviewer} (rapport + pertanyaan), satu
    perekam/pengamat (catatan, \emph{tally}, pengawas keselamatan). Jangan keduanya
    menghadap subjek, dua lawan satu mengintimidasi.
  \item \textbf{\emph{Debrief} hari yang sama,} rekonsiliasi tiga sudut bukti (angka audit,
    perilaku teramati, pandangan dinyatakan) sebelum siapa pun menulis kesimpulan.
\end{itemize}

% =====================================================================
\chapter{Penampilan Lapangan per Zona}
Prinsip: \textbf{cocokkan register, jangan melampauinya}. Di situs bergradien kelas, satu
busana tidak akan bekerja. Berlebihan di sumbu SCBD membuat pekerja menggolongkan tim
sebagai ``asuransi/MLM/sales''; berlebihan di Karet Tengsin membuat tim terbaca sebagai
otoritas, orang luar, atau pembeli yang mengincar lahan. Target: ``kebiasaan yang
dipelajari'', tampak wajar di zona itu, satu tingkat ke arah register subjek, tidak pernah
di atasnya.

\smallskip\noindent\textbf{Blend vs Badge.} \emph{Blend} untuk observasi/\emph{tally}:
profil rendah, sesuai zona, tanpa \emph{clipboard} keluar, ponsel diskret, agar orang
berperilaku normal. \emph{Badge} untuk \emph{intercept}: tampilkan identitas mahasiswa
(\emph{lanyard} dan ID akademi, kaos polos, \emph{clipboard} berisi lembar persetujuan dan
\emph{tally}). Aturan: \emph{blend} saat mengamati, \emph{badge} saat bertanya.

\smallskip\noindent\textbf{Anti-seragam.} Jangan tampak penegak (PKL takut Satpol PP,
patroli sekitar 17.00--21.00 di Benhil): hindari rompi seragam, gaya ``taktis'' serba
hitam, ID-\emph{lanyard} bergaya pemerintah, atau memegang kamera seperti surveyor. Jangan
tampak kru riset komersial. Jangan mencolok kaya di panggung belakang (simpan gawai premium;
papan asumsi sendiri menandai \emph{copet}/\emph{rawan}).

\begin{longtable}{@{}p{2.8cm}p{2.8cm}p{4.0cm}p{3.4cm}@{}}
\toprule
\textbf{Zona} & \textbf{Register} & \textbf{Pakai} & \textbf{Hindari} \\
\midrule
\endhead
\textbf{Sumbu/SCBD} (kantoran, transit) & \emph{Junior-office smart-casual} & Kemeja kerah rapi atau kaos polos bersih, chino/jins bersih, sepatu tertutup, \emph{lanyard} terlihat & Celana pendek, kaos slogan, tampak sales, tampak kuyup/berantakan \\
\addlinespace
\textbf{Simpul transit} (Dukuh Atas, Jl.\ Blora) & Kasual-praktis & Sama tetapi lebih kasual, mudah bergerak di kerumunan & Apa pun yang memperlambat di keramaian \\
\addlinespace
\textbf{Seam Benhil} (PKL) & Polos kasual, ramah & Kaos sederhana, sopan, \emph{beli sesuatu dulu} & Tampilan resmi/audit; postur kamera ke depan \\
\addlinespace
\textbf{Panggung belakang} (Karet Tengsin) & Berpakaian sederhana & Pakaian sederhana tertutup, rendah hati; duduk dan pesan di warung & Gawai mencolok, ponsel mahal di tangan, bahasa tubuh ``memeriksa'' \\
\addlinespace
\textbf{CFD Minggu} (rekreasi) & Sporty-kasual & \emph{Athleisure} menyatu sempurna & Busana kantor (tidak pas) \\
\bottomrule
\end{longtable}

\smallskip\noindent\textbf{Cuaca dan praktis (untuk sekitar 33\textdegree C dan kemungkinan
hujan).} Bahan ringan cepat-kering; sepatu jalan tertutup nyaman; payung lipat/ponco
ringan (Jumat berisiko hujan 25\%, akhir pekan lebih); tabir surya dan air; sopan menutup
bahu dan lutut, terutama ke panggung belakang dan dengan lansia.

\smallskip\noindent\textbf{Anggota perempuan dan pimpinan Sintani (malam/sensitif).}
Praktis, sopan, \emph{alas kaki mobile} (mungkin perlu meninggalkan situasi cepat); ponsel
mudah dijangkau. Ini keselamatan lapangan pragmatis, bukan menyalahkan korban: tim mengkaji
koridor tempat 4 dari 5 perempuan melaporkan pelecehan ruang publik. Pegang aturan
\textbf{berpasangan}; mencocokkan gender pewawancara dengan responden perempuan menurunkan
ancaman yang dirasa.

\smallskip\noindent\textbf{Daftar bawaan.} Terlihat: \emph{lanyard}/ID akademi, \emph{clipboard}
berisi persetujuan + \emph{tally} + kartu pertanyaan, pena. Diskret: ponsel (yang sederhana
di panggung belakang); perekam hanya setelah persetujuan terpisah. Praktis: air,
payung/ponco, tabir surya, uang kecil (untuk membeli dari PKL/warung yang diwawancara).

% =====================================================================
\chapter{Wawancara Ahli, Etika, dan Logistik}

\section{Wawancara Ahli (Jalur Terpisah)}
Berbeda dari \emph{intercept} jalanan, ini terjadwal, semi-terstruktur, berorientasi
hambatan (kesenjangan Bab~II: ``apa yang \emph{mencegah} fasilitas pejalan kaki yang
nyaman''):
\begin{itemize}
  \item \textbf{Anies Baswedan} (menunggu konfirmasi ADA): revitalisasi 2017--2022, apa
    yang menghambat \emph{walkability} yang nyaman, logika relokasi PKL, apa yang tidak
    dapat diselesaikan kebijakan.
  \item \textbf{Karsa City Lab / Urun Daya Kota:} rantai mobilitas perempuan, pembingkaian
    ruang publik yang aman dan nyaman.
  \item \textbf{KRPA / organisasi terafiliasi:} data dan metode pelecehan ruang publik,
    untuk menginformasikan (bukan menggiring) \emph{intercept} perempuan.
\end{itemize}
Perlakukan klaim ahli sebagai konteks untuk diuji di lapangan, bukan sebagai kesimpulan.

\section{Logistik Hari Lapangan}
\textbf{Tiga simpul jangkar:} (1) Dukuh Atas/Stasiun Sudirman untuk klaster komuter +
transit + ojol, (2) \emph{seam} Benhil untuk PKL + kontras depan/belakang, (3) satu gang
Karet Tengsin untuk panggung belakang. Jalankan satu \emph{pass} puncak pagi dan satu
\emph{pass} sore$\rightarrow$malam pada segmen yang \emph{sama} untuk kontras siang/malam.
Tambahkan sesi CFD Minggu hanya sebagai kasus kontras.

\smallskip\noindent\textbf{Perjalanan.} Tim mengkaji simpul transit yang akan dilalui
sendiri. Datang dengan KRL ke Stasiun Sudirman atau MRT ke Dukuh Atas dan \emph{catat
jalan \emph{last-mile} sendiri sebagai data pilot}. Hindari menyetir; parkir adalah ekonomi
informal kampung, bukan kemudahan tim. Siapkan untuk panas tengah hari: bawa air, jadwalkan
audit 11.00--14.00 di sisi teduh.

\smallskip\noindent\textbf{Tempat \emph{debrief}.} Warteg Benhil atau kafe Dukuh Atas
berfungsi sebagai titik \emph{debrief} murah di dalam \emph{fieldsite}; warung kampung
sekaligus pintu masuk riset dan meja \emph{debrief}. \emph{Debrief} hari yang sama selagi
ingatan segar, rekonsiliasi tiga sudut bukti sebelum menulis kesimpulan.

\section{Kartu Etika (Bawa Selalu)}
\begin{enumerate}
  \item Persetujuan lisan sebelum setiap \emph{intercept}; persetujuan terpisah untuk merekam.
  \item \emph{Intercept} sensitif dan malam dipimpin Sintani, sistem berpasangan.
  \item Kaburkan wajah; hormati ketentuan WeatherKit/Google/OSM.
  \item Pertanyaan netral dan dua arah; asumsi pada lembar terpisah, tidak pernah di dalam pertanyaan.
  \item PKL dan warga kampung rentan, pahami mereka, jangan mengaudit.
  \item Siapa pun yang tidak nyaman, hentikan dan ucapkan terima kasih.
  \item Tidak ada temuan ditulis sebelum data terkumpul.
\end{enumerate}

\bigskip
\noindent\rule{\textwidth}{0.4pt}\\
{\footnotesize Sumber yang melandasi profil ini (jejak Bab~II/metodologi): survei pelecehan
ruang publik KRPA 2022; Kompas Megapolitan tentang tumpukan ojol Stasiun Sudirman (Mei 2025);
Mojok dan Kumparan tentang Karet Tengsin di belakang Sudirman Park (2026/2025); Tempo, Media
Indonesia, dan Detik tentang pembersihan PKL Benhil dan patroli Satpol PP (2018, berlanjut);
beragam laporan 2025--2026 tentang ritme komuter Dukuh Atas dan Car Free Day; Knight Frank via
\emph{Expat Life in Indonesia} (submarket kantor Sudirman); Afkara \& Kusuma (2020) tentang norma
jarak berjalan Jakarta; ITDP Indonesia (2019) dan Dwisadana \& Widjajanti (2024) seperti pada
tinjauan pustaka. Kerangka wawancara: Cialdini; Spradley (1979); Fisher \& Geiselman; Miller \&
Rollnick; Tickle-Degnen \& Rosenthal (1990); Tourangeau \& Yan (2007).\par}

\end{document}
